\documentclass[11pt, a4paper]{article}
\usepackage[utf8]{inputenc}
\usepackage[T1]{fontenc}
\usepackage[ngerman]{babel}
\usepackage{geometry}
\usepackage{helvet}
\usepackage{xcolor}
\usepackage{colortbl}
\usepackage{tabularx}
\usepackage{parskip}

% Seitenränder anpassen
\geometry{left=2cm, right=2cm, top=2cm, bottom=2cm}

% Schriftartwechsel auf Sans-Serif
\renewcommand{\familydefault}{\sfdefault}

% Farben definieren
\definecolor{hmblue}{RGB}{0, 51, 102} % HM-Blau
\definecolor{white}{RGB}{255, 255, 255}

\begin{document}

\section*{Projektsteckbrief}
\textbf{Thema:} Generative 'Architecture is Code': KI-basierte Synthese modularer SysML v2 Modelle aus heterogenen Bestandsdaten

\vspace{0.5cm}

\begin{tabularx}{\textwidth}{|>{\bfseries}l|X|}
\hline
\rowcolor{hmblue} \textcolor{white}{Bereich} & \textcolor{white}{Inhalt} \\ \hline
1. Ausgangslage & 
Im Brownfield-Engineering liegen Systeminformationen unstrukturiert und fragmentiert vor (CAD, Excel, PDF). Die manuelle Überführung in MBSE ist fehleranfällig und führt zu monolithischen Modellen ('Silo-Problem'). Implizites Wissen geht bei der reinen Datenkonvertierung verloren. \\ \hline
2. Zielsetzung & 
Entwicklung eines KI-gestützten 'AiC-Generators' unter Nutzung von SysML v2. Das Ziel ist die automatisierte Erzeugung valider, modularer Systemarchitekturen aus Altdaten. Fehlende Schnittstellen sollen durch LLM-basierte Mustererkennung antizipiert werden. \\ \hline
3. Vorgehen & 
\textbf{1. Analyse:} State-of-the-Art Recherche zu SysML v2 und Generative Engineering. \newline
\textbf{2. Konzeption:} Entwurf der dreistufigen Pipeline (Ingestion, Pattern Matching, Synthesis). \newline
\textbf{3. Implementierung:} Entwicklung des Prototypen. \newline
\textbf{4. Validierung:} Durchführung von Blind-Tests und Integration isolierter Teilmodelle. \\ \hline
4. Nutzen & 
\textbullet\ Automatisierte Digitalisierung von Bestandsanlagen ('Schuhschachtel' zu Modell). \newline
\textbullet\ Vermeidung von 'Dead-End'-Architekturen durch intelligente Antizipation. \newline
\textbullet\ Sicherstellung der Modularität durch standardisierten SysML v2 Code. \\ \hline
5. Zeitraum & 
15.03.2026 bis 15.09.2026 (6 Monate) \\ \hline
6. Beteiligte & 
\textbf{Bearbeiter:} Moritz Diez \newline
\textbf{Betreuer HM:} [Name einfügen] \newline
\textbf{Firmenbetreuer:} Kai Heinrich (R + S)\\ \hline
\end{tabularx}

\vspace{1cm}
\textbf{Meilensteinplan (Grob):}
\begin{itemize}
    \item \textbf{M1 (15.04.):} Konzeptreview und Anforderungsdefinition abgeschlossen.
    \item \textbf{M2 (31.05.):} Prototyp des Generators funktionsfähig implementiert.
    \item \textbf{M3 (15.07.):} Validierung durch Fallstudien erfolgreich durchgeführt.
    \item \textbf{M4 (15.09.):} Abgabe der Masterarbeit.
\end{itemize}

\end{document}